\documentclass[10pt,a4paper]{amsart}
\usepackage{amsmath,amssymb,amsthm,graphicx}
\usepackage{verbatim}
\newenvironment{reference}{\comment}{\endcomment}
%\newenvironment{reference}{}{}
\newenvironment{slogan}{\comment}{\endcomment}
\newenvironment{history}{\comment}{\endcomment}
% For commutative diagrams we use Xy-pic
\usepackage[all]{xy}

% We use 2cell for 2-commutative diagrams.
\xyoption{2cell}
% Theorem environments.
%
\theoremstyle{plain}
\newtheorem{theorem}[subsection]{Theorem}
\newtheorem{proposition}[subsection]{Proposition}
\newtheorem{lemma}[subsection]{Lemma}

\theoremstyle{definition}
\newtheorem{definition}[subsection]{Definition}
\newtheorem{example}[subsection]{Example}
\newtheorem{exercise}[subsection]{Exercise}
\newtheorem{situation}[subsection]{Situation}

\theoremstyle{remark}
\newtheorem{remark}[subsection]{Remark}
\newtheorem{remarks}[subsection]{Remarks}

\numberwithin{equation}{subsection}

% Macros
%
\def\lim{\mathop{\mathrm{lim}}\nolimits}
\def\colim{\mathop{\mathrm{colim}}\nolimits}
\def\Spec{\mathop{\mathrm{Spec}}}
\def\Hom{\mathop{\mathrm{Hom}}\nolimits}
\def\Ext{\mathop{\mathrm{Ext}}\nolimits}
\def\SheafHom{\mathop{\mathcal{H}\!\mathit{om}}\nolimits}
\def\SheafExt{\mathop{\mathcal{E}\!\mathit{xt}}\nolimits}
\def\Sch{\mathit{Sch}}
\def\Mor{\mathop{\mathrm{Mor}}\nolimits}
\def\Ob{\mathop{\mathrm{Ob}}\nolimits}
\def\Sh{\mathop{\mathit{Sh}}\nolimits}
\def\NL{\mathop{N\!L}\nolimits}
\def\CH{\mathop{\mathrm{CH}}\nolimits}
\def\proetale{{pro\text{-}\acute{e}tale}}
\def\etale{{\acute{e}tale}}
\def\QCoh{\mathit{QCoh}}
\def\Ker{\mathop{\mathrm{Ker}}}
\def\Im{\mathop{\mathrm{Im}}}
\def\Coker{\mathop{\mathrm{Coker}}}
\def\Coim{\mathop{\mathrm{Coim}}}

% Boxtimes
%
\DeclareMathSymbol{\boxtimes}{\mathbin}{AMSa}{"02}

%
% Macros for moduli stacks/spaces
%
\def\QCohstack{\mathcal{QC}\!\mathit{oh}}
\def\Cohstack{\mathcal{C}\!\mathit{oh}}
\def\Spacesstack{\mathcal{S}\!\mathit{paces}}
\def\Quotfunctor{\mathrm{Quot}}
\def\Hilbfunctor{\mathrm{Hilb}}
\def\Curvesstack{\mathcal{C}\!\mathit{urves}}
\def\Polarizedstack{\mathcal{P}\!\mathit{olarized}}
\def\Complexesstack{\mathcal{C}\!\mathit{omplexes}}
% \Pic is the operator that assigns to X its picard group, usage \Pic(X)
% \Picardstack_{X/B} denotes the Picard stack of X over B
% \Picardfunctor_{X/B} denotes the Picard functor of X over B
\def\Pic{\mathop{\mathrm{Pic}}\nolimits}
\def\Picardstack{\mathcal{P}\!\mathit{ic}}
\def\Picardfunctor{\mathrm{Pic}}
\def\Deformationcategory{\mathcal{D}\!\mathit{ef}}
\usepackage{cleveref}
\crefname{originallemma}{original lemma}{original lemmas}
\Crefname{originallemma}{Original lemma}{Original lemmas}
\makeatletter
\newlabel{lemma-chow-ring-projective-bundle}{{02TX}{source}{}{sourcepointer-1}{}}
\newlabel{lemma-chow-ring-projective-bundle@cref}{{[originallemma][0][]02TX}{[1][0][]source}{}{}{}}
\makeatother
\crefname{originalremark}{original remark}{original remarks}
\Crefname{originalremark}{Original remark}{Original remarks}
\makeatletter
\newlabel{remark-pullback-cohomology}{{0B7F}{source}{}{sourcepointer-2}{}}
\newlabel{remark-pullback-cohomology@cref}{{[originalremark][0][]0B7F}{[1][0][]source}{}{}{}}
\makeatother
\crefname{originalreference}{original reference}{original references}
\Crefname{originalreference}{Original reference}{Original references}
\makeatletter
\newlabel{situation-setup}{{02QL}{source}{}{sourcepointer-3}{}}
\newlabel{situation-setup@cref}{{[originalreference][0][]02QL}{[1][0][]source}{}{}{}}
\makeatother
\title{The splitting principle: a projective flat modification making finitely many bundles filtrations by invertible sheaves}
\author{The Stacks Project Authors}
\date{stacks-project:chow.tex:local-snapshot-20261012; GNU FDL 1.2 or later}
\begin{document}
\maketitle
\noindent\textit{Editorial note.} All delivered mathematical text is the author's own, copied from the pinned Stacks Project chapter file chow.tex; no statement or proof step was rewritten, reordered or removed.\par\medskip
\noindent\textit{Editorial note.} Preamble provenance: the pinned chapter file opens with \textbackslash{}input\{preamble\} and carries no preamble text of its own, so the project's single Stacks-wide preamble staged at sources/stacks-chow/preamble.tex is excerpted verbatim as three spans, lines 8-12 (the verbatim package together with the project's own reference, slogan and history display-environment definitions, including the author's commented-out older form of the reference environment), lines 14-18 (Xy-pic setup) and lines 154-217 (the native theorem declarations and mathematical macros). The lines 8-12 span is required because a counted statement or its surrounding text may use those display environments, which the project typesets as comments. Nothing was added, removed or redefined; the wrapper class and the amsmath/amssymb/amsthm/graphicx base packages are supplied by the assembler.\par\medskip
\noindent\textit{Editorial note.} Same-chapter pointers are located at their \textbackslash{}label lines inside the pinned chapter file; cross-chapter pointers are located in the project's own tag list staged in the snapshot (sources/stacks-chow/tags.txt), at the verbatim lines carrying their canonical tags.\par\medskip
\noindent\textit{Editorial note.} The author's closing sentence at source line 8463, 'Some details omitted.', is retained verbatim at the end of the counted proof; the induction, the projective bundle step and the composition are all inside the span, and only routine verifications are left to the reader.\par\medskip
\noindent\textit{Allowed standard prerequisites (editorial).}\begin{itemize}
\item Injectivity of pullback along a projective bundle, and the existence of an invertible quotient of the pulled back bundle. Same chapter chow.tex, label lemma-chow-ring-projective-bundle, invoked at proof line 8458. Gives that the pullback along the projective bundle of one of the bundles is injective on Chow groups after any base change, which is the induction bookkeeping, and produces the short exact sequence exhibiting an invertible quotient.
\item The remark identifying the induced map on bivariant groups with restriction. Same chapter chow.tex, label remark-pullback-cohomology, invoked at proof line 8442. Translates injectivity on Chow groups after base change into injectivity of the restriction map on the bivariant theory, which is the additional statement of the lemma.
\end{itemize}\medskip
\ifdefined\hypertarget\hypertarget{sourcepointer-1}{}\fi
\noindent\textit{Source reference.} Original lemma 02TX: sources/stacks-chow/pinned-original-source.tex, source lines 6932--6932. Navigation pointer for \textbackslash{}ref\{lemma-chow-ring-projective-bundle\} used inside the counted span (line 8458); the argument that consumes it is retained in the span.\par\medskip
\ifdefined\hypertarget\hypertarget{sourcepointer-2}{}\fi
\noindent\textit{Source reference.} Original remark 0B7F: sources/stacks-chow/pinned-original-source.tex, source lines 6289--6289. Navigation pointer for \textbackslash{}ref\{remark-pullback-cohomology\} used inside the counted span (line 8442); the argument that consumes it is retained in the span.\par\medskip
\ifdefined\hypertarget\hypertarget{sourcepointer-3}{}\fi
\noindent\textit{Source reference.} Original reference 02QL: sources/stacks-chow/pinned-original-source.tex, source lines 1372--1372. Navigation pointer for \textbackslash{}ref\{situation-setup\} used inside the counted span (line 8430); the argument that consumes it is retained in the span.\par\medskip
\begin{lemma}
\label{lemma-splitting-principle}
Let $(S, \delta)$ be as in Situation \ref{situation-setup}. Let $X$ be locally
of finite type over $S$. Let $\mathcal{E}_i$ be a finite collection of
locally free $\mathcal{O}_X$-modules of rank $r_i$. There exists a projective
flat morphism $\pi : P \to X$ of relative dimension $d$ such that
\begin{enumerate}
\item for any morphism $f : Y \to X$ the map
$\pi_Y^* : \CH_*(Y) \to \CH_{* + d}(Y \times_X P)$ is injective, and
\item each $\pi^*\mathcal{E}_i$ has a filtration
whose successive quotients $\mathcal{L}_{i, 1}, \ldots, \mathcal{L}_{i, r_i}$
are invertible ${\mathcal O}_P$-modules.
\end{enumerate}
Moreover, when (1) holds the restriction map $A^*(X) \to A^*(P)$
(Remark \ref{remark-pullback-cohomology}) is injective.
\end{lemma}

\begin{proof}
We may assume $r_i \geq 1$ for all $i$. We will prove the lemma by induction
on $\sum (r_i - 1)$. If this integer is $0$, then $\mathcal{E}_i$
is invertible for all $i$ and we conclude by taking $\pi = \text{id}_X$.
If not, then we can pick an $i$ such that $r_i > 1$ and consider the
morphism $\pi_i : P_i = \mathbf{P}(\mathcal{E}_i) \to X$.
We have a short exact sequence
$$
0 \to \mathcal{F} \to \pi_i^*\mathcal{E}_i \to \mathcal{O}_{P_i}(1) \to 0
$$
of finite locally free $\mathcal{O}_{P_i}$-modules of ranks $r_i - 1$,
$r_i$, and $1$. Observe that $\pi_i^*$ is injective on chow groups
after any base change by the projective bundle formula
(Lemma \ref{lemma-chow-ring-projective-bundle}).
By the induction hypothesis applied to the finite locally free
$\mathcal{O}_{P_i}$-modules $\mathcal{F}$ and $\pi_{i'}^*\mathcal{E}_{i'}$
for $i' \not = i$, we find a morphism $\pi : P \to P_i$ with
properties stated as in the lemma. Then the composition
$\pi_i \circ \pi : P \to X$ does the job. Some details omitted.
\end{proof}

\end{document}
